\section{Simultaneous chromatic bases and the filtration}
\label{sec:spectral}

The coefficient theorem computes an algebraic continuous-cochain
complex, whereas the desired conclusion concerns specified maps of
spectra. We compare that complex with the ordinary height-\(t\)
residue of the completed Morava descent resolution. The same
sphere-cochain representatives then lift every coefficient basis
vector through the resulting totalization tower, forcing the actual
product maps to be \(K(t)\)-equivalences. The three steps below
construct the common maps, identify their residue coefficients and
constant action, and prove that their images form a basis. Fix
\(p>n+1\), put \(S=S_p^\wedge\), and write
\[
D=L_{K(n)}S.
\]
Choose a finite field \(k\) containing \(\mathbb F_{p^a}\) for
\(1\leq a\leq n\), with \(f=[k:\Fp]\) prime to \(p\). For example,
\(f=\operatorname{lcm}(1,\ldots,n)\) has these properties. Write
\[
\Gamma=\Gal(k/\Fp),\qquad
G_n=P_n\rtimes\Gamma,
\]
where the action on the ordinary height-\(n\) stabilizer factors through
\(\Gal(\mathbb F_{p^n}/\Fp)\). Let \(E_n=E_n(k)\) be Morava
\(E\)-theory with this field of constants. Its coefficient ring is
\[
\pi_0E_n=W(k)[[x_1,\ldots,x_{n-1}]],
\]
with its maximal-ideal topology. The extension of constants allows the
same source resolution to be used for every height test below.

\subsection{Completed descent and exact functors}

All the tensor products in the following Amitsur resolution are
initially taken in the \(K(n)\)-local category. Set
\begin{equation}
\mathcal A^q=
L_{K(n)}\bigl(E_n^{\wedge(q+1)}\bigr),\qquad q\geq0.
\label{spec:amitsur}
\end{equation}
The ordinary smash product inside the parentheses is followed by the
indicated localization. Completed Morava cooperations identify this
cosimplicial spectrum with
\[
\mathcal A^q\simeq\operatorname{Map}^c(G_n^q,E_n),
\]
where the right side uses maximal-ideal-completed functions. We use
standard inhomogeneous coordinates: the first coface is the
semilinear stabilizer action, and the remaining cofaces are the
group-nerve maps. The augmentation is the unit of \(D\).
For the standard constant field \(k_0=\mathbb F_{p^n}\), these
completed cooperations have the structured realization of
Devinatz--Hopkins. Their evaluation composites (2.3), (2.5), and
(2.7), formed from the action and multiplication, define the natural
cohomology isomorphism (2.6), with the completed coefficient Hopf
algebroid described by Proposition~2.2
\cite[Proposition~2.2 and (2.3), (2.5)--(2.7)]{DevinatzHopkins}.
The weakly contractible mapping components and rectification in their
Theorems~4.1 and~3.2 and Proposition~4.6 make these maps coherent.
Their Proposition~4.10 and Construction~4.11 construct diagrams over
finite continuous \(G_n\)-sets, with the full simplex category,
all cofaces, and all codegeneracies
\cite[Theorems~3.2 and~4.1, Propositions~4.6 and~4.10, and
Construction~4.11]{DevinatzHopkins}. Their Theorem~1(iii) identifies
the augmented object with \(D\).

Here \(\operatorname{Map}^c\) denotes our chosen completed
function-spectrum model for this rectified diagram, functorial for
continuous pullback in its profinite parameters. In the printed
coordinates of Devinatz--Hopkins, the inverse action occurs in the
last coface \cite[Definition~4.18 and Lemma~4.22]{DevinatzHopkins}.
On the structured model, parameter pullback and the rectified action
map realize the following change from a standard inhomogeneous
\(q\)-cochain \(c\) to their coordinates:
\[
(T_qc)(g_1,\ldots,g_q)=
g_1^{-1}c(g_1g_2^{-1},\ldots,g_{q-1}g_q^{-1},g_q),
\qquad T_0=\mathrm{id}.
\]
The inverse at \((a_1,\ldots,a_q)\) evaluates at
\((a_1\cdots a_q,a_2\cdots a_q,\ldots,a_q)\) and acts on the value by
\(a_1\cdots a_q\). These are inverse maps of the completed function
spectra by the structured action law. Applying the same substitutions
to the rectified structure maps intertwines all cofaces and
codegeneracies; the unital multiplicative action preserves the unit
and multiplication. Thus this transport supplies the stated
inhomogeneous convention at the spectrum level.
We explain the extension to the chosen \(k\). Put
\(Q=\operatorname{Gal}(k/k_0)\). The unramified coefficient extension
\(E_n(k_0)\to E_n(k)\) is finite free, and its Galois comparison
\[
E_n(k)\otimes_{E_n(k_0)}E_n(k)
 \xrightarrow{\;\sim\;}\prod_{\sigma\in Q}E_n(k)
\]
is an equivalence in the \(K(n)\)-local module category: on
coefficients it is the finite unramified identity
\(W(k)\otimes_{W(k_0)}W(k)\cong\prod_Q W(k)\), and finite freeness
removes higher Tor. Base-changing the standard cooperation formula
on both ends and using this identity indexes its factors by the
lifts in
\[
P_n\rtimes\operatorname{Gal}(k/\Fp)
 \longrightarrow P_n\rtimes\operatorname{Gal}(k_0/\Fp).
\]
This gives the displayed formula for \(G_n\); iteration gives all
higher terms. Naturality for the unit, multiplication, and rectified
action transports \(T_\bullet\) to these factors. In our inhomogeneous
coordinates the first coface therefore remains the semilinear action,
with the remaining nerve maps unchanged. The same
finite Galois comparison identifies descent through \(Q\) with the
standard constant-field theory, so the augmentation still has
source \(D\).

\begin{proposition}[Exact tests of completed descent]
\label{spec:exact-descent}
The augmentation \(D\to\operatorname{Tot}\mathcal A^\bullet\) is an
equivalence. Moreover, if \(F\) is an exact functor from spectra to a
stable \(\infty\)-category admitting countable limits, then
\begin{equation}
F(D)\xrightarrow{\;\simeq\;}
\operatorname{Tot}\bigl(F(\mathcal A^\bullet)\bigr).
\label{spec:exact-totalization}
\end{equation}
The transformed partial-totalization tower is quickly convergent.
For spectrum-valued \(F\), its totalization spectral sequence has
this convergence control. In particular, if its
\(E_2\)-page lies in a bounded range of cohomological degrees, it
converges strongly with a finite filtration in each total degree.
\end{proposition}

\begin{proof}
Morava \(E\)-theory attached to the chosen perfect field \(k\) is
descendable in the \(K(n)\)-local category
\cite[Proposition~10.10]{Mathew}. Descendability says that the augmented
partial-totalization tower has nilpotent error; equivalently, it
belongs to the thick class generated by constant and nilpotent
towers \cite[Propositions~3.10,~3.20 and~3.27]{Mathew}. Here a nilpotent
tower is one for which a fixed finite number of consecutive transition
maps has null composite. Such a tower has zero inverse limit.
This is the quickly convergent form of descent, whose universal
totalization property is explained in
\cite[Definition~2.19 and Propositions~2.20,~2.21 and~2.26]
{MathewNilpotence}.

The inclusion of \(K(n)\)-local spectra into spectra is exact.
We verify that each full partial totalization here is a finite limit.
For \(s\geq0\), send a nonempty subset of \(\{0,\ldots,s\}\) to its
ordered set in \(\Delta_{\leq s}\). This functor is homotopy initial.
Indeed, its comma category at \([q]\), \(q\leq s\), is the nonempty
face poset of the complex with simplices
\[
(i_0,j_0),\ldots,(i_b,j_b),\qquad
0\leq i_0<\cdots<i_b\leq s,\quad
0\leq j_0\leq\cdots\leq j_b\leq q.
\]
Call this complex \(P(s,q)\). It is contractible for \(s\geq q\).
Start with the cone with vertex \((s,q)\) on \(P(s-1,q)\), and
attach in descending order the cones with vertices \((s,j)\),
\(j<q\). The attaching subcomplex is \(P(s-1,j)\), which is
contractible by induction since \(j\leq s-1\). The first cone
is contractible even when \(q=s\); the induction begins with
\(P(0,0)\), a point. This proves the claim.

Consequently the full partial totalization is the limit of the
punctured \((s+1)\)-cube with vertex
\(L_{K(n)}(E_n^{\wedge|T|})\) at nonempty
\(T\subseteq\{0,\ldots,s\}\) and unit-insertion maps.
This is the finite cubical comparison in standard descent machinery;
compare \cite[Proposition~2.14]{MNNDescent}. It applies to the full
Amitsur totalization, including its codegeneracies. Thus the inclusion
and then \(F\) preserve these finite limits, the augmentation, and
the stated nilpotence of the error tower. The inverse limit of the transformed
error tower is therefore zero. This proves
\eqref{spec:exact-totalization} and carries along the convergence
control. Bounded cohomological support then gives the stated finite,
complete filtration on the abutment.
\end{proof}

Thus an ordinary residue smash can be applied after
\eqref{spec:amitsur}. The terms to which it is applied remain the
completed terms of that resolution.

\subsection{Common spherical classes}

For a profinite group \(G\), use the continuous sphere-cochain spectrum
\(C^*_{\mathrm{cts}}(G;S)\) of
Proposition~\ref{prop:modification-transport}. One may construct its
cosimplicial terms by locally constant approximation with finite
\(p\)-power coefficients and finite Postnikov truncations, followed
by the separated Postnikov and \(p\)-complete limits. The unit
\(S\to E_n\) gives a natural map of cosimplicial spectra
\begin{equation}
C^\bullet_{\mathrm{cts}}(G_n;S)
\longrightarrow\operatorname{Map}^c(G_n^\bullet,E_n).
\label{spec:constant-map}
\end{equation}
Indeed, it gives the map on locally constant functions at every
finite coefficient stage. Completeness of the target extends it to
the indicated limits. At finite Postnikov range this is exactly the
continuous cochain construction used to obtain the spherical classes
in Theorem~\ref{thm:stable-generators}; passing to the Postnikov limit
preserves this identification. The parameter pullbacks defining
\(T_\bullet\) commute with these constant-section maps, and its value
action commutes with the unit because every action map is unital.
Thus \eqref{spec:constant-map} uses the same transported structured
cosimplicial model, preserving all cofaces and codegeneracies, the unit,
and multiplication through the indicated limits.

Write $\alpha^P_{n,i}\in H^{2i-1}(P_n,\Zp)$ for the transport
of the matrix class $\alpha_{n,i}$ under
Proposition~\ref{prop:modification-transport}, and write
$\alpha^P_{n,I}=\prod_{i\in I}\alpha^P_{n,i}$.
Its reduction modulo $p$, extended to $k$, is the class $\xi_I$
used in Proposition~\ref{prop:constant-module}.

\Needspace{10\baselineskip}
\begin{proposition}[Simultaneous representatives]
\label{spec:global-classes}
There are classes
\[
y_i\in\pi_{1-2i}D,\qquad 1\leq i\leq n,
\]
represented by constant sphere-cochain classes whose integral
Hurewicz images on \(P_n\) are \(\alpha^P_{n,i}\). For
\(I=\{i_1<\cdots<i_r\}\), put
\[
y_I=y_{i_1}\cdots y_{i_r},\qquad y_\varnothing=1.
\]
These are represented by actual sphere-cochain classes with
Hurewicz image \(\alpha^P_{n,i_1}\cdots\alpha^P_{n,i_r}\). The class
\(y_\varnothing\) is the unit \(S\to D\).
\end{proposition}

\begin{proof}
Theorem~\ref{thm:stable-generators} and
Proposition~\ref{prop:modification-transport} give classes
\[
u_i\in\pi_{1-2i}C^*_{\mathrm{cts}}(P_n;S)
\]
with these integral Hurewicz images. The range applies because
\(n\leq p-2\).

By Remark~\ref{rem:const-galois}, the action of $\Gamma$ on
$H^*_{\mathrm{cts}}(P_n;\Zp)$ is trivial.

As \(f\) is a unit in \(\Zp\), restriction from \(G_n\) identifies
continuous sphere cochains with the \(\Gamma\)-homotopy fixed
points of the \(P_n\)-cochains, and the averaging idempotent realizes
invariants on homotopy groups. In concrete terms, normalized transfer
lifts
\[
\overline u_i=\frac1f\sum_{\gamma\in\Gamma}\gamma u_i
\]
to a class on \(G_n\). Its Hurewicz image on \(P_n\) remains
\(\alpha^P_{n,i}\). This also follows from the finite-group
homotopy-fixed-point spectral sequence: all its positive
\(\Gamma\)-cohomology vanishes, since multiplication by \(f\) is
invertible on its coefficient groups.

Map these classes through the totalization of
\eqref{spec:constant-map} and use
Proposition~\ref{spec:exact-descent} to obtain the \(y_i\).
Multiplication supplies the stated representatives of all the
\(y_I\). The empty product maps to the augmentation unit itself.
\end{proof}

\subsection{Ordinary residue coefficients}

Fix \(1\leq t<n\), and put \(m=n-t\). Let \(E_t=E_t(k)\) be
height-\(t\) Morava \(E\)-theory and let \(\kappa_t\) be its
iterated residue module, obtained by taking cofibers of the regular
sequence
\[
p,u_1,\ldots,u_{t-1}.
\]
Put \(I_t=(p,u_1,\ldots,u_{t-1})\subseteq\pi_0E_t\).
Regularity identifies its homotopy groups, as graded
\(\pi_*E_t\)-modules, with the quotient scalar field
\[
\pi_*\kappa_t\cong\kappa_{t,*}:=\pi_*E_t/I_t
=k[v^{\pm1}],\qquad |v|=2.
\]
This residue has the Bousfield class of \(K(t)\); consequently its
ordinary smash detects \(K(t)\)-equivalences. The periodicity here
belongs to the test theory.

The scalar action on the residue tower comes from \(E_t\).
For each displayed generator \(x\) of \(I_t\), let \(E_t/x\) be its
cofiber as an \(E_t\)-module. Evenness and the fact that \(x\) is a
nonzerodivisor on \(\pi_*E_t\) give \(\pi_1(E_t/x)=0\).
Applying \([-,E_t/x]_{E_t}\) to the cofiber triangle gives an injection
\[
[E_t/x,E_t/x]_{E_t}\hookrightarrow
[E_t,E_t/x]_{E_t}=\pi_0(E_t/x).
\]
The map \(x\,\mathrm{id}_{E_t/x}\) restricts to zero in the group on
the right, so it is nullhomotopic. The iterated residue is the relative
tensor product of these individual cofibers. Tensoring the chosen
nullhomotopy on the \(x\)-factor shows that \(x\) acts nullhomotopically
on \(\kappa_t\). Smashing that homotopy with a spectrum is natural in
the spectrum, so it also gives a nullhomotopy on the residue
cosimplicial diagram and on all its partial totalizations and fibers.
Consequently their homotopy groups are modules over the quotient
\(\kappa_{t,*}\), compatibly with the tower filtration. Each scalar
may be represented by a lift in \(\pi_*E_t\); its action is a map of
the tower, and different lifts give the same action on homotopy.

Retain the algebraic marking tower of
Proposition~\ref{prop:frame-torsor} and the cochain convention of
Definition~\ref{rings:bounded-cochains}: with
\[
A=k[[x_t,\ldots,x_{n-1}]],\qquad B=A[1/x_t],
\]
write
\begin{equation}
\mathcal B_{n,t}
 =B_{\mathrm{conn},\GL_m(\Zp)}
 =\colim_K B_{K,\GL_m(\Zp)}.
\label{spec:connected-coefficients}
\end{equation}
Here \(K\) runs through open connected-marking levels. The
\(\GL_m(\Zp)\) subscript means that the integral \'etale
Tate module is unmarked. In particular, no additional general linear
group parameter is part of \(\mathcal B_{n,t}\).

We first record the flatness that is needed before taking a residue.

\begin{lemma}[Flat continuous function modules]
\label{spec:flat-functions}
Let \(R=W(k)[[x_1,\ldots,x_{n-1}]]\), with maximal ideal
\(\mathfrak m\), and let \(Q\) be a profinite set. The module
\(N_Q=C_{\mathrm{cts}}(Q,R)\) is \(\mathfrak m\)-adically
complete and flat over \(R\). Its reductions are
\[
N_Q/\mathfrak m^aN_Q
 =C_{\mathrm{lc}}(Q,R/\mathfrak m^a).
\]
Quotienting by \(p,x_1,\ldots,x_{t-1}\) gives
\(C_{\mathrm{cts}}(Q,A)\).
\end{lemma}

\begin{proof}
A function into the finite ring \(R/\mathfrak m^a\) is constant on
some finite clopen partition of \(Q\); lifting its finitely many
values proves surjectivity of reduction. The kernel equality follows
by continuous coefficient division among the finitely many monomial
generators of \(\mathfrak m^a\). This division is performed in the
power-series coordinates, including the \(p\)-adic coefficient
expansion, and is continuous for their adic topologies. The same
coordinate argument applies to the displayed height ideal.

The locally constant function module modulo \(\mathfrak m^a\) is
a filtered colimit of finite free \(R/\mathfrak m^a\)-modules,
indexed by finite clopen partitions, and is therefore flat.
Furthermore
\[
N_Q=\varprojlim_a C_{\mathrm{lc}}(Q,R/\mathfrak m^a).
\]
The adic flatness criterion over the Noetherian ring \(R\) now gives
flatness of \(N_Q\). Equivalently, lifting a basis of
\(C_{\mathrm{lc}}(Q,k)\) and successively lifting modulo
\(\mathfrak m^a\) identifies \(N_Q\) with a completed free
\(R\)-module. This also proves its asserted completeness.
\end{proof}

\begin{lemma}[Ordinary residue of the completed cobar]
\label{lem:ordinary-residue}
For the completed terms \(\mathcal A^\bullet\) in
\eqref{spec:amitsur}, followed by ordinary smash with \(\kappa_t\),
there is a natural quasi-isomorphism of cochain complexes
\begin{equation}
N^*\pi_r(\kappa_t\wedge\mathcal A^\bullet)
\simeq
\kappa_{t,r}\otimes_k
 C^*_{\mathrm{cts}}(P_n;\mathcal B_{n,t}).
\label{spec:residue-complex}
\end{equation}
Here $N^*$ denotes normalization of a cosimplicial abelian group.
The right side uses one finite algebraic connected-marking stage and
one bound on the pole order on each compact cochain domain. It is zero
when \(r\) is odd. The comparison carries the unit to \(1\) and the
constant-cochain action to the actual constant cup action of
Proposition~\ref{prop:constant-module}.
\end{lemma}

\begin{proof}
\emph{Ordinary cooperations before residue.}
For \(M=\operatorname{Map}^c(Q,E_n)\), the ordinary module identity is
\[
E_t\wedge M\simeq
(E_t\wedge E_n)\otimes_{E_n}M.
\]
Lemma~\ref{spec:flat-functions} gives flatness of \(\pi_*M\) over
\(E_{n,*}\), so the module Kunneth spectral sequence has no positive
\(\operatorname{Tor}\). Its convergence is bounded by the finite
graded global dimension of the regular ring \(E_{n,*}\). Thus
\begin{equation}
\pi_*(E_t\wedge M)=
\pi_*(E_t\wedge E_n)\otimes_{E_{n,*}}
 C_{\mathrm{cts}}(Q,E_{n,*}).
\label{spec:ordinary-kunneth}
\end{equation}
Ordinary Landweber exact base change identifies the first factor
with the two-sided base change of the formal-group-isomorphism
Hopf algebroid. This is the ordinary cooperation algebra; the
Landweber flatness assertions on both sides apply to it
\cite[Example~0.1(d), Lemma~2.2 and Corollary~2.3]
{HoveyStrickland}.
One may compute using \(BP_*BP\): isomorphisms between the typical
laws are the formal-group sums of their typical curves, and the
ratio of the periodic frames supplies the invertible linear
coefficient. Thus this computation includes all isomorphisms of the
underlying formal groups
\cite[Definition~2.40, Remark~2.41(3) and Corollary~2.45]
{GoerssFormalGroups}.

\emph{The height ideal and the algebraic marking ring.}
Flatness and \eqref{spec:ordinary-kunneth} keep the test-side
sequence \(p,u_1,\ldots,u_{t-1}\) regular. Indeed, it is regular
on the cooperation algebra, and tensoring its successive injections
with the flat function module preserves injectivity. Taking its
iterated cofibers therefore produces only even homotopy groups.
Under the universal formal-group isomorphism the successive height
ideals agree on the two sides; their next generators differ by a
unit modulo the preceding ideal. Consequently the residue imposes
\[
p=x_1=\cdots=x_{t-1}=0,
\]
and invertibility of the test height-\(t\) coefficient inverts
\(x_t\). Lemma~\ref{spec:flat-functions} then leaves
\(C_{\mathrm{cts}}(Q,A)[1/x_t]\) on the source side.

We make the remaining cooperation algebra explicit, including the
source constants and periodicity. The two constant fields are
initially independent, and
\[
k\otimes_{\Fp}k\cong\prod_{\sigma\in\Gamma}k,
\qquad a\otimes b\longmapsto(a\sigma(b))_\sigma,
\]
where the first factor is the test field. Let
\(\mathcal B_{n,t}^{[\sigma]}\) be the connected-marking algebra
after the source constants are embedded into the test field by
\(\sigma\). Its universal marking is an isomorphism from the Honda
height-\(t\) group to that source deformation. In particular,
\(\mathcal B_{n,t}^{[1]}=\mathcal B_{n,t}\). With the test period
\(v\) fixed, the ordinary cooperation calculation gives an
isomorphism of graded algebras
\begin{equation}
\label{spec:residue-cooperation-algebra}
\frac{\pi_*(E_t\wedge E_n)}
     {(p,u_1,\ldots,u_{t-1})}
\ \cong\
\prod_{\sigma\in\Gamma}
 \bigl(\kappa_{t,*}\otimes_k\mathcal B_{n,t}^{[\sigma]}\bigr).
\end{equation}
Here the ideal on the left acts through the test factor.

To check the map and its inverse, use degree-zero formal-group
coordinates and let \(\alpha\) denote the Honda-to-source
isomorphism. If \(v_n\) is the source period, its linear coefficient
is \(c=v_n/v\). Comparison of the leading terms in
\([p]_{\mathrm{source}}\alpha=\alpha[p]_{\mathrm{Honda}}\)
gives
\[
x_t c^{p^t-1}=1.
\]
Thus the cooperation data determine the embedding \(\sigma\) and
the full marking \(\alpha\), including \(c\). Conversely, a marking
and the fixed test period determine \(v_n=cv\); rescaling the
coordinates by these periods makes \(\alpha\) strict and recovers
the corresponding \(BP_*BP\)-isomorphism. These constructions are
inverse and commute with composition. In particular, the source
period is already recorded by the linear coefficient of the marking;
there is no further frame parameter. No \'etale Tate basis is chosen.

The full algebraic isomorphism ring is a filtered union
of finite \'etale marking algebras
\cite[Lecture~14, Theorem~1]{LurieFormalGroups}.
These stages retain coefficients of full markings together with their
extension equations, as in Proposition~\ref{prop:frame-torsor}; they
are not schemes of arbitrary finite-bud isomorphisms.

\emph{The topology produced by the tensor product.}
For a finite marking algebra \(J\) over \(B\), choose an idempotent
presentation \(J=eB^a\). Then
\begin{equation}
J\otimes_A C_{\mathrm{cts}}(Q,A)
=e\bigl(C_{\mathrm{cts}}(Q,A)[1/x_t]\bigr)^a.
\label{spec:bounded-functions}
\end{equation}
The right side consists exactly of functions admitting a common
power of \(x_t\) as denominator and continuous coordinates in a
finite \(A\)-lattice. For example, \(J\cap A^a\) is such a
lattice: it is finite over the Noetherian ring \(A\), it spans
\(J\) after inverting \(x_t\), and it is closed in \(A^a\).
Multiplying by a common denominator places any element of the right
side in continuous functions into this lattice; the converse follows
from its inclusion in \(A^a\). Passing to the algebraic union of
finite \(J\)'s requires one finite stage for each element of the
tensor product. This proves exactly the bounded-pole, common-stage
convention on every profinite parameter \(Q\). No completion is
taken after inverting \(x_t\) or forming this algebraic union.

\emph{Cofaces and constants.}
Apply this calculation with \(Q=G_n^q\), retaining every
source-embedding factor in
\eqref{spec:residue-cooperation-algebra}. For the first coface,
composition with the deformation-change isomorphism transports the
tautological marking by the actual stabilizer action. Every finite
coefficient of this composite uses only finitely many coefficients
of the two isomorphisms. The continuous action preserves the height
ideal and takes \(x_t\) to a unit times \(x_t\) modulo that ideal.
Compactness of the stabilizer parameter therefore gives one finite
marking stage and one pole bound for the resulting coefficient.
The other cofaces and the codegeneracies are pullbacks along the
usual profinite group-nerve maps. They preserve the same convention.
The unit is the constant function \(1\); multiplication is ordinary
multiplication of the marking coefficients. This proves compatibility
with constant cup products as well as with the differential.

It remains to take source-field descent. With the
test field fixed, \(\Gamma\) permutes the source-embedding factors
simply transitively; the full product of the corresponding
\(P_n\)-cochain coefficients is a coinduced module for
\(P_n\subset G_n\). Group descent, or Shapiro's lemma, identifies
its \(G_n\)-cochains with the \(P_n\)-cochains of the
identity-embedding factor. This is the step that changes the
termwise cooperation calculation into the quasi-isomorphism of
\(P_n\)-cochain complexes in \eqref{spec:residue-complex}.
Equivalently, exact \(\Gamma\)-descent
on the full product followed by evaluation at that factor gives this
identification. Taking invariants of an isolated factor would not
make sense, since that factor is not \(\Gamma\)-stable. The test
periodicity is fixed throughout, so \(\kappa_{t,r}\) is an
untwisted coefficient line. This gives
\eqref{spec:residue-complex}, naturally for constants and their
products.
\end{proof}

The use of ordinary smash here also respects the scalar convention of
the theorem. The cofiber of \(S_{(p)}\to S_p^\wedge\) is rational:
it is \(p\)-local and multiplication by \(p\) on it is an
equivalence, since the map is an equivalence modulo \(p\).
Its smash with \(\kappa_t\) is zero. Consequently
\(\kappa_t\wedge S\simeq\kappa_t\), and for an \(S\)-module
\(M\), ordinary residue smash agrees with the corresponding
relative \(S\)-module test. All the maps used below are thus maps
of \(S\)-modules.

\subsection{The specified bases survive}

\begin{theorem}[Simultaneous height bases]
\label{thm:height-bases}
The common classes of Proposition~\ref{spec:global-classes} satisfy,
for every \(1\leq t<n\),
\begin{equation}
L_{K(t)}\left(
\bigvee_{I\subseteq\{1,\ldots,n-t\}}
 \Sigma^{-d(I)}S\xrightarrow{(y_I)}D
\right)
\quad\text{is an equivalence},
\label{spec:height-basis}
\end{equation}
where \(d(I)=\sum_{i\in I}(2i-1)\), and the empty component is the
unit.
\end{theorem}

\begin{proof}
Use Proposition~\ref{spec:exact-descent} with the exact functor
\(\kappa_t\wedge-\). Lemma~\ref{lem:ordinary-residue} identifies
the totalization spectral sequence as
\begin{equation}
E_2^{s,r}=\kappa_{t,r}\otimes_k
 H^s_{\mathrm{cts}}(P_n;\mathcal B_{n,t})
\Longrightarrow\pi_{r-s}(\kappa_t\wedge D).
\label{spec:test-spectral-sequence}
\end{equation}
By Proposition~\ref{prop:constant-module}, this page is free over
$\kappa_{t,*}$ with basis $\xi_I\cdot1$, for
$I\subseteq\{1,\ldots,m\}$ and $m=n-t$. These are the images
of the reductions of the transported integral classes
$\alpha^P_{n,I}$.
The basis vector indexed by \(I\) lies in bidegree
\((s,r)=(d(I),0)\). In particular, the page is zero for
\(s<0\) or \(s>m^2\). This statement uses the constant cup action
and the nonzero image of \(1\) in that proposition.

We now identify the actual representative of each basis vector.
There are two filtrations to distinguish: the source tower of
constant sphere cochains is built from connective spectra, whereas
the residue tower computing \(\kappa_t\wedge D\) is periodic.
We first locate a representative in the source filtration and then
transport it to the residue filtration.

The empty subset is represented by the unit itself. For a nonempty
subset put \(b=d(I)>0\). The sphere-cochain representative of
\(y_I\) has filtration at least \(b\). Indeed, the
\((b-1)\)-st partial totalization of the connective sphere-cochain
diagram is \((1-b)\)-connective and hence has zero homotopy in
degree \(-b\). The representative therefore lifts to the fiber
of the map to that partial totalization. Its component in filtration
\(b\), modulo the cochain boundary, is its integral Hurewicz
class, whose restriction to \(P_n\) is \(\alpha^P_{n,I}\).
The map \eqref{spec:constant-map} sends this chosen lift to the
completed descent tower. Ordinary residue smash is exact and
preserves the partial totalizations and their fibers, so it transports
the lift to the residue tower as well. Lemma~\ref{lem:ordinary-residue}
identifies its leading component in
\eqref{spec:test-spectral-sequence} with precisely the specified
image of \(\alpha^P_{n,I}\). Connectivity was needed only in the
source tower.

These representatives come from the source totalization and hence
supply compatible lifts through every subsequent partial
totalization. Their leading terms cannot support an outgoing
differential. Products give the representatives for all subsets. For
a coefficient in \(\kappa_{t,*}\), choose a lift in \(\pi_*E_t\).
Its natural action on the residue tower carries these compatible lifts
to representatives of the corresponding multiple of each basis vector.
Therefore every element of the \(E_2\)-page has such a lift, so
\(d_2=0\). Inductively, if the earlier differentials vanish, the
same representatives span the next page and force its outgoing
differential to vanish. This proves that all differentials vanish;
in particular none of the specified vectors can be lost as an
incoming boundary.

Strong convergence and the bound \(0\leq s\leq m^2\) now give a
finite filtration on \(\pi_*(\kappa_t\wedge D)\). The actual
images of the \(y_I\) lift its specified associated-graded basis.
The coefficient action preserves every filtration step, so this is a
filtration by graded \(\kappa_{t,*}\)-submodules.
Finite filtered linear algebra over the graded field
\(\kappa_{t,*}\) shows that these images themselves form a basis:
independence and spanning follow successively from their lowest
nonzero filtration components. It follows that the actual map
\[
\bigvee_{I\subseteq\{1,\ldots,m\}}
 \Sigma^{-d(I)}\kappa_t
\longrightarrow\kappa_t\wedge D
\]
is an equivalence. Since \(\kappa_t\) detects
\(K(t)\)-equivalences and ordinary residue smash agrees with the
relative scalar test, this proves \eqref{spec:height-basis} with all
its stated components.
\end{proof}

\subsection{Rational dimensions and completion of the proof}

The last block of the filtration is rational. We use the following
independent calculation of the rational local sphere.

\begin{theorem}[Barthel--Schlank--Stapleton--Weinstein]
\label{thm:rational-sphere}
For every prime \(p\) and height \(n\geq1\), the rational homotopy
of the \(K(n)\)-local sphere is the exterior \(\Qp\)-algebra
\[
\pi_*L_{K(n)}S\otimes_{\Z}\Q
\cong\bigwedge_{\Qp}(\zeta_1,\ldots,\zeta_n),
\qquad |\zeta_i|=1-2i,
\]
with its natural scalar action induced by the unit. In particular,
\begin{equation}
\dim_{\Qp}\pi_{-b}L_0D
 =\#\{I\subseteq\{1,\ldots,n\}:d(I)=b\}
\qquad(b\in\Z).
\label{spec:rational-dimensions}
\end{equation}
\end{theorem}

\begin{proof}
This is \cite[Theorem~A]{BSSW}, with the generator degrees displayed
in homotopical grading. Replacing the sphere by its \(p\)-completion
does not change its positive-height localization. The dimension
formula is the monomial basis of the exterior algebra. Only this
additive formula and the natural scalar action are needed below; no
identification of the individual \(\zeta_i\) with the classes
\(y_i\) is required.
\end{proof}

\begin{proof}[Proof of Theorem~\ref{thm:main}]
Proposition~\ref{spec:global-classes} supplies the one family
\(y_i\), its ordered products, and \(y_\varnothing=1\).
For \(n>1\), take \(z_I=y_I\), for
\(I\subseteq\{1,\ldots,n-1\}\), and take the specified map
\(u=y_\varnothing:S\to D\) to be the unit.
Theorem~\ref{thm:height-bases} verifies
all the simultaneous local basis hypotheses of
Proposition~\ref{prop:assembly};
Theorem~\ref{thm:rational-sphere} verifies its rational dimension
hypothesis. That proposition constructs the successive quotient
blocks and pulls them back to a filtration of \(L_{n-1}D\), with
the required individual cofibers. Under the first-stage identification
\(F_1=L_{n-1}S\), the composite \(F_1\to L_{n-1}D=X_{n,p}\) is the
localized unit. All maps and cofibers are \(S\)-linear, as required.

When \(n=1\), Proposition~\ref{spec:global-classes} still supplies
\(y_1\); only the positive-height tests are vacuous.
The unit calculation in Section~\ref{sec:assembly} and the dimensions
of Theorem~\ref{thm:rational-sphere} verify the remaining hypotheses
of Proposition~\ref{prop:assembly}. Its height-one case gives the
two cofibers \(H\Qp\) and \(\Sigma^{-1}H\Qp\), with the composite
from the first stage to \(X_{1,p}\) equal to the canonical localization unit.
\end{proof}
